\documentclass[11pt,a4paper]{article}
\usepackage[utf8]{inputenc}
\usepackage[spanish,es-tabla]{babel}
\usepackage{amsmath,amsfonts,amssymb}
\usepackage{geometry}
\usepackage{graphicx}
\usepackage{xcolor}

\usepackage{enumitem}
\usepackage{tikz}

\geometry{margin=2.5cm}

\title{\textbf{Soluci\'on Examen Final 2023\\TAF STAR\\IMT Atlantique}}
\author{}
\date{}

\begin{document}

\maketitle

\section*{Ejercicio 1: Training Sequence and Channel Code}

\textbf{Contexto del problema:}

Sistema de comunicaci\'on con:
\begin{itemize}
    \item Modulaci\'on: QPSK
    \item Rapidez de s\'imbolo: $\frac{1}{T} = 2$ MHz = $2 \times 10^6$ sym/s
    \item Bloques de $N$ s\'imbolos
    \item Tiempo de coherencia del canal: $T_c = 32 \times 10^{-6}$ s = 32 μs
    \item Ecualizaci\'on necesaria (canal con multipaso)
    \item Secuencia de entrenamiento (training sequence) para estimar el canal
\end{itemize}

\textbf{Funcionamiento:}
\begin{enumerate}
    \item Se env\'ia un bloque de $N$ s\'imbolos
    \item Se estima el canal usando la training sequence
    \item Se ecualiza el bloque recibido
    \item Se env\'ia el siguiente bloque
\end{enumerate}

\subsection*{Pregunta 1: Elegir tama\~no del bloque $N$}

\textbf{Enunciado:} Choose the size of the transmission block (i.e. the value of $N$) and justify your choice.

\vspace{0.3cm}
\textbf{Soluci\'on paso a paso:}

\textbf{Paso 1: Comprender el problema}

El canal tiene un tiempo de coherencia $T_c = 32$ μs, lo que significa que el canal permanece aproximadamente constante durante este tiempo. Despu\'es, cambia significativamente.

\textbf{Pregunta clave:} ¿Cu\'antos s\'imbolos se pueden transmitir antes de que el canal cambie demasiado?

\textbf{Paso 2: Calcular el n\'umero m\'aximo de s\'imbolos por bloque}

Periodo s\'imbolo:
\[
T_{sym} = \frac{1}{2 \times 10^6} = 0.5 \times 10^{-6} \text{ s} = 0.5 \text{ μs}
\]

N\'umero de s\'imbolos que caben en el tiempo de coherencia:
\[
N_{max} = \frac{T_c}{T_{sym}} = \frac{32 \times 10^{-6}}{0.5 \times 10^{-6}} = 64 \text{ s\'imbolos}
\]

\textbf{Paso 3: Consideraciones pr\'acticas}

Si transmitimos m\'as de 64 s\'imbolos:
\begin{itemize}
    \item El canal habr\'a cambiado significativamente
    \item La estimaci\'on inicial ya no es v\'alida
    \item La ecualizaci\'on ser\'a incorrecta
    \item Aumentar\'a la tasa de error
\end{itemize}

Si transmitimos menos de 64 s\'imbolos:
\begin{itemize}
    \item Infrautilizamos el tiempo de coherencia
    \item Necesitamos re-entrenar m\'as frecuentemente
    \item Menor eficiencia (m\'as overhead de training)
\end{itemize}

\textbf{Paso 4: An\'alisis de la Figure 1 (MSE vs Training Sequence Size)}

De la figura proporcionada, observamos que:
\begin{itemize}
    \item Para training sequence peque\~no ($< 20$ s\'imbolos): MSE alto (mala estimaci\'on)
    \item Para training sequence $\approx 20$ s\'imbolos: MSE comienza a estabilizarse
    \item Para training sequence $> 30$ s\'imbolos: MSE ya es bajo y se satura
\end{itemize}

\textbf{Conclusi\'on:} Necesitamos al menos 20-30 s\'imbolos de training para buena estimaci\'on.

\textbf{Paso 5: Estructura del bloque}

Un bloque completo consiste en:
\[
N = N_{training} + N_{data}
\]

Restricciones:
\begin{enumerate}
    \item $N \leq 64$ (l\'imite del tiempo de coherencia)
    \item $N_{training} \geq 20$ (buena estimaci\'on seg\'un Figure 1)
\end{enumerate}

\textbf{Opci\'on \'optima:}

Si elegimos $N_{training} = 20$ s\'imbolos:
\[
N_{data,max} = 64 - 20 = 44 \text{ s\'imbolos de datos}
\]

Si elegimos $N_{training} = 30$ s\'imbolos (m\'as robusto):
\[
N_{data,max} = 64 - 30 = 34 \text{ s\'imbolos de datos}
\]

\textbf{Decisi\'on:}

\[
\boxed{N = 64 \text{ s\'imbolos totales}}
\]

con estructura:
\[
\boxed{N_{training} = 20-30 \text{ s\'imbolos, } N_{data} = 34-44 \text{ s\'imbolos}}
\]

\textbf{Justificaci\'on completa:}

\begin{enumerate}
    \item \textbf{L\'imite f\'isico:} $N = 64$ es el m\'aximo que permite el tiempo de coherencia
    
    \item \textbf{Eficiencia:} Maximizamos el uso del tiempo de coherencia
    
    \item \textbf{Calidad de estimaci\'on:} 20-30 s\'imbolos de training dan MSE suficientemente bajo
    
    \item \textbf{Compromiso:} Balance entre overhead de training (30-46\%) y calidad de estimaci\'on
\end{enumerate}

\textbf{Representaci\'on temporal:}

\begin{center}
\begin{tikzpicture}[scale=0.8]
% Bloque 1
\draw[thick, fill=red!20] (0,0) rectangle (1.5,1.5);
\draw[thick, fill=blue!20] (1.5,0) rectangle (4,1.5);
\node at (0.75,0.75) {\small Train};
\node at (2.75,0.75) {\small Datos};

% Bloque 2
\draw[thick, fill=red!20] (4.5,0) rectangle (6,1.5);
\draw[thick, fill=blue!20] (6,0) rectangle (8.5,1.5);
\node at (5.25,0.75) {\small Train};
\node at (7.25,0.75) {\small Datos};

% Tiempo de coherencia
\draw[<->, thick, green!70!black] (0,-0.5) -- (4,-0.5);
\node[green!70!black, below] at (2,-0.5) {$T_c = 32$ μs};

% N\'umero de s\'imbolos
\node[above] at (2,1.5) {$N = 64$ s\'imbolos};

\end{tikzpicture}
\end{center}

\textbf{Eficiencia del sistema:}

Con $N_{training} = 20$, $N_{data} = 44$:
\[
\eta_{training} = \frac{N_{data}}{N} = \frac{44}{64} = 68.75\%
\]

El 31.25\% del tiempo se usa para training, pero es necesario para mantener la calidad de transmisi\'on.



\subsection*{Pregunta 2: D\'ebito binario \'util sin c\'odigo (BER = $10^{-3}$)}

\textbf{Enunciado:} Compute the useful binary rate (number of information bit/sec) if the needed quality of the transmission is equal to $10^{-3}$.

\vspace{0.3cm}
\textbf{Soluci\'on paso a paso:}

\textbf{Paso 1: Analizar la Figure 2 (Equalizer performance vs MSE)}

De la Figure 2, necesitamos encontrar qu\'e MSE permite lograr BER = $10^{-3}$.

T\'ipicamente, para BER = $10^{-3}$, se requiere:
\[
\text{MSE} \leq 0.01 \text{ (valor t\'ipico, ajustar seg\'un la gr\'afica real)}
\]

\textbf{Paso 2: Determinar el tama\~no de training necesario}

De la Figure 1 (MSE vs Training size):

Para MSE = 0.01:
\[
N_{training} \approx 20 \text{ s\'imbolos}
\]

Por tanto:
\[
N_{data} = 64 - 20 = 44 \text{ s\'imbolos de datos}
\]

\textbf{Paso 3: Calcular bits por s\'imbolo QPSK}

QPSK transmite:
\[
\log_2(4) = 2 \text{ bits por s\'imbolo}
\]

\textbf{Paso 4: Calcular bits de datos por bloque}

\[
\text{Bits de datos por bloque} = N_{data} \times 2 = 44 \times 2 = 88 \text{ bits}
\]

\textbf{Paso 5: Calcular duraci\'on de un bloque}

\[
T_{bloque} = N \times T_{sym} = 64 \times 0.5 \times 10^{-6} = 32 \times 10^{-6} \text{ s} = 32 \text{ μs}
\]

\textbf{Paso 6: Calcular d\'ebito binario \'util}

\[
D_b = \frac{\text{Bits de datos por bloque}}{T_{bloque}} = \frac{88}{32 \times 10^{-6}}
\]

\[
D_b = \frac{88}{32} \times 10^{6} = 2.75 \times 10^{6} \text{ bits/s}
\]

\[
\boxed{D_b = 2.75 \text{ Mbits/s}}
\]

\textbf{Verificaci\'on alternativa:}

D\'ebito de s\'imbolos: $D_s = 2 \times 10^6$ sym/s

Bits por s\'imbolo (QPSK): 2 bits/sym

D\'ebito bruto: $D_{bruto} = 2 \times 2 \times 10^6 = 4$ Mbits/s

Eficiencia por training:
\[
\eta = \frac{44}{64} = 68.75\%
\]

D\'ebito \'util:
\[
D_b = 4 \times 0.6875 = 2.75 \text{ Mbits/s} \quad $\checkmark$
\]

\textbf{Resumen del c\'alculo:}

\begin{center}
\begin{tabular}{|l|r|}
\hline
\textbf{Par\'ametro} & \textbf{Valor} \\
\hline
Rapidez s\'imbolos & 2 Msym/s \\
Modulaci\'on & QPSK (2 bits/sym) \\
D\'ebito bruto & 4 Mbits/s \\
S\'imbolos totales/bloque & 64 \\
S\'imbolos training & 20 \\
S\'imbolos datos & 44 \\
Eficiencia & 68.75\% \\
\textbf{D\'ebito \'util} & \textbf{2.75 Mbits/s} \\
\hline
\end{tabular}
\end{center}

\textbf{Interpretaci\'on:}

El overhead del 31.25\% debido al training es el precio a pagar para:
\begin{itemize}
    \item Mantener BER = $10^{-3}$ (calidad aceptable)
    \item Seguir las variaciones del canal
    \item Ecualizar correctamente el canal con multipaso
\end{itemize}

Sin training, la BER ser\'ia inaceptablemente alta.



\subsection*{Pregunta 3: D\'ebito con c\'odigo de tasa $\frac{1}{2}$ (BER = $10^{-3}$)}

\textbf{Enunciado:} Compute the useful binary rate if we use a channel code of rate $\frac{1}{2}$ and the needed quality is still $10^{-3}$.

\vspace{0.3cm}
\textbf{Soluci\'on paso a paso:}

\textbf{Paso 1: Comprender el c\'odigo de tasa $\frac{1}{2}$}

Un c\'odigo corrector de errores de rendimiento (rate) $R_c = \frac{1}{2}$ significa:

\[
\text{Por cada 1 bit de informaci\'on, se transmiten 2 bits codificados}
\]

O equivalentemente:
\[
\text{Bits de informaci\'on} = \text{Bits codificados} \times \frac{1}{2}
\]

\textbf{Paso 2: Analizar la Figure 3 (Decoder performance)}

La Figure 3 muestra la relaci\'on entre:
\begin{itemize}
    \item Eje X: BER a la entrada del decodificador (antes de correcci\'on)
    \item Eje Y: BER a la salida del decodificador (despu\'es de correcci\'on)
\end{itemize}

\textbf{Lectura de la gr\'afica:}

Para obtener BER salida = $10^{-3}$, se requiere:
\[
\text{BER entrada} \approx 10^{-2} \text{ (valor t\'ipico, ajustar seg\'un gr\'afica real)}
\]

\textbf{Interpretaci\'on:} El c\'odigo puede tolerar 10 veces m\'as errores a su entrada y aun as\'i cumplir el objetivo de BER = $10^{-3}$ a la salida.

\textbf{Paso 3: Determinar nuevo tama\~no de training}

Con BER entrada = $10^{-2}$ (menos exigente que $10^{-3}$):

De la Figure 2, necesitamos menor calidad de ecualizaci\'on:
\[
\text{MSE requerido} \approx 0.1 \text{ (menos estricto)}
\]

De la Figure 1, con MSE = 0.1:
\[
N_{training,nuevo} \approx 10 \text{ s\'imbolos (menos que antes)}
\]

\textbf{Paso 4: Calcular nuevos s\'imbolos de datos}

\[
N_{data,nuevo} = 64 - 10 = 54 \text{ s\'imbolos de datos}
\]

\textbf{Paso 5: Calcular bits codificados por bloque}

\[
\text{Bits codificados} = N_{data,nuevo} \times 2 = 54 \times 2 = 108 \text{ bits}
\]

\textbf{Paso 6: Aplicar el rendimiento del c\'odigo}

\[
\text{Bits de informaci\'on} = \text{Bits codificados} \times R_c = 108 \times \frac{1}{2} = 54 \text{ bits}
\]

\textbf{Paso 7: Calcular d\'ebito binario \'util}

Duraci\'on del bloque (igual que antes):
\[
T_{bloque} = 32 \text{ μs}
\]

D\'ebito:
\[
D_b = \frac{54}{32 \times 10^{-6}} = \frac{54}{32} \times 10^{6} = 1.6875 \times 10^{6} \text{ bits/s}
\]

\[
\boxed{D_b = 1.6875 \text{ Mbits/s} \approx 1.69 \text{ Mbits/s}}
\]

\textbf{Comparaci\'on con caso sin c\'odigo:}

\begin{center}
\begin{tabular}{|l|c|c|}
\hline
\textbf{Par\'ametro} & \textbf{Sin c\'odigo} & \textbf{Con c\'odigo $R=\frac{1}{2}$} \\
\hline
BER objetivo & $10^{-3}$ & $10^{-3}$ \\
BER antes decoder & $10^{-3}$ & $10^{-2}$ \\
MSE requerido & 0.01 & 0.1 \\
Training s\'imbolos & 20 & 10 \\
Datos s\'imbolos & 44 & 54 \\
Bits codificados & 88 & 108 \\
Bits informaci\'on & 88 & 54 \\
Overhead training & 31.25\% & 15.6\% \\
Overhead c\'odigo & 0\% & 50\% \\
\textbf{D\'ebito \'util} & \textbf{2.75 Mbits/s} & \textbf{1.69 Mbits/s} \\
\hline
\end{tabular}
\end{center}

\textbf{An\'alisis:}

\textbf{Reducci\'on de d\'ebito:}
\[
\frac{1.69}{2.75} = 0.614 = 61.4\% \text{ del d\'ebito sin c\'odigo}
\]

\textbf{¿Por qu\'e menos d\'ebito?}

Dos factores:
\begin{enumerate}
    \item \textbf{Overhead del c\'odigo:} 50\% (por cada 1 bit de info, 2 bits transmitidos)
    \item \textbf{Eficiencia mejorada por training:} 84.4\% vs 68.75\%
\end{enumerate}

El factor (1) reduce m\'as de lo que (2) compensa, resultando en menor d\'ebito neto.



\subsection*{Pregunta 4: ¿Vale la pena usar el c\'odigo?}

\textbf{Enunciado:} Is it worth to use this channel code?

\vspace{0.3cm}
\textbf{Soluci\'on paso a paso:}

\textbf{An\'alisis cuantitativo:}

\begin{center}
\begin{tabular}{|l|c|c|}
\hline
\textbf{M\'etrica} & \textbf{Sin c\'odigo} & \textbf{Con c\'odigo} \\
\hline
D\'ebito \'util & 2.75 Mbits/s & 1.69 Mbits/s \\
BER & $10^{-3}$ & $10^{-3}$ \\
Training overhead & 31.25\% & 15.6\% \\
Robustez & Baja & Alta \\
Complejidad & Baja & Media \\
\hline
\end{tabular}
\end{center}

\textbf{Ventajas del c\'odigo:}

\begin{enumerate}
    \item \textbf{Mayor robustez:}
    \begin{itemize}
        \item Tolera BER entrada = $10^{-2}$ vs $10^{-3}$
        \item $10\times$ m\'as errores tolerables
        \item Menos sensible a variaciones del canal
    \end{itemize}
    
    \item \textbf{Menor overhead de training:}
    \begin{itemize}
        \item Solo 10 s\'imbolos vs 20 s\'imbolos
        \item Menos tiempo dedicado a estimaci\'on
        \item Re-estimaci\'on m\'as frecuente posible si necesario
    \end{itemize}
    
    \item \textbf{Margen de operaci\'on:}
    \begin{itemize}
        \item Puede funcionar con menor SNR
        \item M\'as tolerante a interferencias
        \item Mejor en condiciones adversas
    \end{itemize}
\end{enumerate}

\textbf{Desventajas del c\'odigo:}

\begin{enumerate}
    \item \textbf{Menor d\'ebito:}
    \begin{itemize}
        \item 1.69 vs 2.75 Mbits/s
        \item Reducci\'on del 38.5\%
    \end{itemize}
    
    \item \textbf{Mayor complejidad:}
    \begin{itemize}
        \item Codificador en emisor
        \item Decodificador en receptor (Viterbi, LDPC, Turbo...)
        \item Mayor latencia de procesamiento
        \item Mayor consumo energ\'etico
    \end{itemize}
    
    \item \textbf{Mayor latencia:}
    \begin{itemize}
        \item Retardo del decodificador
        \item Puede ser cr\'itico en aplicaciones tiempo-real
    \end{itemize}
\end{enumerate}

\textbf{Decisi\'on:}

\[
\boxed{\text{DEPENDE de la aplicaci\'on y condiciones del canal}}
\]

\textbf{Usar c\'odigo corrector SI:}

\begin{itemize}
    \item El canal es muy variable o ruidoso
    \item Se necesita alta fiabilidad (transmisi\'on de datos cr\'iticos)
    \item El d\'ebito no es la prioridad principal
    \item Se puede tolerar mayor latencia
    \item La potencia de transmisi\'on es limitada (c\'odigo permite menor SNR)
    \item Ejemplo: telemetr\'ia satelital, comunicaciones submarinas, IoT con bater\'ia limitada
\end{itemize}

\textbf{NO usar c\'odigo corrector SI:}

\begin{itemize}
    \item Se necesita m\'aximo d\'ebito
    \item El canal es relativamente bueno y estable
    \item Latencia baja es cr\'itica (voz en tiempo real, gaming)
    \item Complejidad/coste deben minimizarse
    \item Se puede aumentar potencia de transmisi\'on si necesario
    \item Ejemplo: streaming de video local, comunicaciones de muy corta distancia
\end{itemize}

\textbf{Recomendaci\'on para este caso:}

Dado que:
\begin{itemize}
    \item El tiempo de coherencia es solo 32 μs (canal muy variable)
    \item Se requiere ecualizaci\'on (multipaso severo)
    \item BER objetivo es $10^{-3}$ (comunicaci\'on de datos, no voz)
\end{itemize}

\[
\boxed{\text{S\'I, vale la pena usar el c\'odigo corrector}}
\]

\textbf{Justificaci\'on:}

El canal es claramente desafiante (coherencia de solo 32 μs). La mayor robustez del c\'odigo compensa la p\'erdida de d\'ebito, especialmente porque:

\begin{enumerate}
    \item Permite operar con menor SNR $\to$ mayor alcance o menor potencia
    
    \item Reduce overhead de training $\to$ m\'as eficiente en uso del tiempo de coherencia
    
    \item Mayor fiabilidad $\to$ menos retransmisiones $\to$ mejor throughput efectivo
\end{enumerate}

\textbf{C\'alculo de throughput efectivo:}

Sin c\'odigo:
\begin{itemize}
    \item D\'ebito: 2.75 Mbits/s
    \item Si BER = $10^{-3}$, PER (packet error) $\approx 10\%$ para paquetes de 100 bits
    \item Throughput efectivo: $2.75 \times 0.9 = 2.475$ Mbits/s (con retransmisiones)
\end{itemize}

Con c\'odigo:
\begin{itemize}
    \item D\'ebito: 1.69 Mbits/s
    \item BER = $10^{-3}$ con m\'as margen, PER $\approx 1\%$
    \item Throughput efectivo: $1.69 \times 0.99 = 1.67$ Mbits/s
\end{itemize}

Aunque hay p\'erdida de d\'ebito (1.67 vs 2.475 Mbits/s), la mayor confiabilidad justifica el uso del c\'odigo en canales dif\'iciles.



\subsection*{Pregunta 5: D\'ebitos con BER = $10^{-4}$}

\textbf{Enunciado:} Compute the useful binary rate with and without the channel code if the needed quality is now $10^{-4}$.

\vspace{0.3cm}
\textbf{Soluci\'on paso a paso:}

\textbf{CASO A: Sin c\'odigo corrector (BER = $10^{-4}$)}

\textbf{Paso 1:} De la Figure 2, para BER = $10^{-4}$:
\[
\text{MSE requerido} \approx 0.001 \text{ (muy estricto)}
\]

\textbf{Paso 2:} De la Figure 1, para MSE = 0.001:
\[
N_{training} \approx 40 \text{ s\'imbolos (mucho training necesario)}
\]

\textbf{Paso 3:} S\'imbolos de datos:
\[
N_{data} = 64 - 40 = 24 \text{ s\'imbolos}
\]

\textbf{Paso 4:} Bits de informaci\'on:
\[
\text{Bits info} = 24 \times 2 = 48 \text{ bits}
\]

\textbf{Paso 5:} D\'ebito:
\[
D_b = \frac{48}{32 \times 10^{-6}} = 1.5 \times 10^{6} \text{ bits/s}
\]

\[
\boxed{D_b^{sin\_codigo} = 1.5 \text{ Mbits/s}}
\]

\textbf{CASO B: Con c\'odigo corrector de tasa $\frac{1}{2}$ (BER = $10^{-4}$)}

\textbf{Paso 1:} De la Figure 3, para BER salida = $10^{-4}$:
\[
\text{BER entrada} \approx 5 \times 10^{-3} \text{ (valor t\'ipico)}
\]

\textbf{Paso 2:} De la Figure 2, para BER = $5 \times 10^{-3}$:
\[
\text{MSE requerido} \approx 0.05
\]

\textbf{Paso 3:} De la Figure 1, para MSE = 0.05:
\[
N_{training} \approx 15 \text{ s\'imbolos}
\]

\textbf{Paso 4:} S\'imbolos de datos:
\[
N_{data} = 64 - 15 = 49 \text{ s\'imbolos}
\]

\textbf{Paso 5:} Bits codificados:
\[
\text{Bits codificados} = 49 \times 2 = 98 \text{ bits}
\]

\textbf{Paso 6:} Bits de informaci\'on (aplicando $R_c = \frac{1}{2}$):
\[
\text{Bits info} = 98 \times \frac{1}{2} = 49 \text{ bits}
\]

\textbf{Paso 7:} D\'ebito:
\[
D_b = \frac{49}{32 \times 10^{-6}} = 1.53125 \times 10^{6} \text{ bits/s}
\]

\[
\boxed{D_b^{con\_codigo} = 1.53 \text{ Mbits/s}}
\]

\textbf{Comparaci\'on completa:}

\begin{center}
\begin{tabular}{|l|c|c|c|c|}
\hline
\textbf{BER objetivo} & \multicolumn{2}{c|}{\textbf{$10^{-3}$}} & \multicolumn{2}{c|}{\textbf{$10^{-4}$}} \\
\hline
\textbf{Configuraci\'on} & Sin c\'odigo & Con c\'odigo & Sin c\'odigo & Con c\'odigo \\
\hline
Training (s\'imbolos) & 20 & 10 & 40 & 15 \\
Datos (s\'imbolos) & 44 & 54 & 24 & 49 \\
Bits codificados & 88 & 108 & 48 & 98 \\
Bits informaci\'on & 88 & 54 & 48 & 49 \\
\textbf{D\'ebito (Mbits/s)} & \textbf{2.75} & \textbf{1.69} & \textbf{1.5} & \textbf{1.53} \\
\hline
\end{tabular}
\end{center}

\textbf{An\'alisis de resultados:}

\textbf{Observaci\'on clave para BER = $10^{-4}$:}

\[
\boxed{D_b^{con\_codigo} > D_b^{sin\_codigo} \quad (1.53 > 1.5 \text{ Mbits/s})}
\]

\textbf{¡El c\'odigo corrector AUMENTA el d\'ebito \'util cuando se requiere BER muy baja!}

\textbf{Explicaci\'on:}

\begin{enumerate}
    \item \textbf{Sin c\'odigo:} Para lograr BER = $10^{-4}$ directamente:
    \begin{itemize}
        \item Se requiere MSE muy bajo (0.001)
        \item Necesita 40 s\'imbolos de training (62.5\% del bloque)
        \item Solo 24 s\'imbolos para datos (37.5\%)
        \item Enorme overhead de training
    \end{itemize}
    
    \item \textbf{Con c\'odigo:} Para lograr BER = $10^{-4}$ a la salida:
    \begin{itemize}
        \item Solo se requiere BER = $5 \times 10^{-3}$ a la entrada del decoder
        \item MSE menos exigente (0.05)
        \item Solo 15 s\'imbolos de training (23.4\%)
        \item 49 s\'imbolos para datos (76.6\%)
        \item Aunque el c\'odigo tiene overhead 50\%, el menor training compensa
    \end{itemize}
\end{enumerate}

\textbf{Balance de overheads:}

Sin c\'odigo:
\[
\text{Overhead total} = 62.5\% \text{ (solo training)}
\]

Con c\'odigo:
\[
\text{Overhead training} = 23.4\%
\]
\[
\text{Overhead c\'odigo} = 50\%
\]
\[
\text{Eficiencia neta} = (1 - 0.234) \times 0.5 = 0.383 = 38.3\%
\]
\[
\text{Overhead total equivalente} = 61.7\%
\]

Aunque los overheads son similares (62.5\% vs 61.7\%), la ventaja del c\'odigo es que:
\begin{itemize}
    \item Es m\'as robusto
    \item Menos sensible a variaciones
    \item Mayor margen de operaci\'on
\end{itemize}

\textbf{Conclusi\'on general:}

\[
\boxed{\text{Para BER muy bajas ($10^{-4}$ o menor), el c\'odigo corrector es ESENCIAL}}
\]

\textbf{Ventaja decisiva:} Permite lograr BER muy bajas sin necesidad de estimaci\'on perfecta del canal, lo cual ser\'ia casi imposible en canales tan variables.

\textbf{Lecci\'on importante:}

C\'odigos correctores no solo mejoran la fiabilidad, sino que en canales dif\'iciles con alta calidad requerida, pueden incluso \textbf{aumentar el throughput \'util} al reducir dram\'aticamente el overhead de training/pilot necesario.




\vspace{1cm}
\begin{center}
\rule{0.8\linewidth}{0.4pt}\\
\textit{Fin del examen}
\end{center}

\end{document}
